\subsection{Parameters, hiding leaves, and authenticated oracle semantics}\label{sec:full:oracles}

Throughout this chapter $\F$ is a finite field of \emph{odd}
characteristic; $N=2^m$, $M=2^{m+s}$, $s\geq1$,
and $\Omega_j$ is the subgroup of size $M_j=M/2^j$
from Section~\ref{sec:kbfoldfold:codewords}.
Fix $1\leq B\leq N$, $B_j=\lceil B/2^j\rceil$,
$N_j=2^{m-j}$ and $D_j=N_j+B_j\leq M_j$.
No characteristic-two trace or Artin--Schreier
identity is invoked in this chapter.

\begin{definition}[Canonical field encoding]\label{def:full:field-encoding}
Fix a public injective map
$\mathsf{enc}_{\F}:\F\hookrightarrow\{0,1\}^b$
with $b=\lceil\log_2|\F|\rceil$.
The image of this map is the set of \emph{valid}
encodings; a relation circuit must check validity,
rather than interpreting every $b$-bit string as
a field element. All field arithmetic in the
relation is arithmetic of $\F$, not arithmetic
of the commitment group.
\end{definition}

\begin{definition}[Hiding commitment to a field symbol]\label{def:full:symbol-commit}
Let $G=\langle g\rangle=\langle h\rangle$
be an odd-prime-order group as in
Definition~\ref{def:realize:group}.
For $a\in\F$ with encoding
$(a_1,\ldots,a_b)$, sample independent
$r_1,\ldots,r_b\leftarrow\mathbb Z_P$
and form
\[
 \mathsf{VC}(a;\mathbf r)
   :=\bigl(g^{a_1}h^{r_1},\ldots,
            g^{a_b}h^{r_b}\bigr)\in G^b.
\]
Each component is a Pedersen bit commitment.
The map $\mathsf{VC}$ is \emph{not}
assumed to be additive over $\F$.
\end{definition}

\begin{lemma}[Perfect uniformity of a committed field symbol]\label{lem:full:leaf-uniform}
For every fixed $a\in\F$ and fresh independent
uniform $\mathbf r$, the public vector
$\mathsf{VC}(a;\mathbf r)$ is uniform on $G^b$.
For any finite collection of possibly correlated
field symbols, their independently blinded
commitment vectors are jointly uniform on the
corresponding power of $G$, independently of
all symbol values.
\end{lemma}
\begin{proof}
Apply Lemma~\ref{lem:realize:hiding}
to each of the $b$ Pedersen bit commitments.
Their blindings are independent, so their
joint law is uniform on $G^b$, for every
fixed encoding. Conditioning on any fixed
collection of field symbols, the entire
collection of commitments is thus a product
uniform law independent of those symbols.
Averaging over any distribution on the
symbols preserves the same product law.
\end{proof}

\begin{definition}[Domain-separated Merkle trees]\label{def:full:merkle}
Fix a deterministic hash
$H:\{0,1\}^*\to\{0,1\}^{\lambda_H}$
and a canonical injective serialization
$\mathsf{ser}$ for group elements, indices and
domain identifiers. For each level $j$ and
ordered index $x\in\Omega_j$, put
\[
 L_{j,x}:=H(\mathsf{ser}(\texttt{leaf},j,x,
                        \mathsf{VC}(Q_j(x);\mathbf r_{j,x}))).
\]
Internal binary nodes hash
$\mathsf{ser}(\texttt{node},j,\text{left},\text{right})$.
Pad non-power-of-two levels with a publicly
fixed convention if needed. Write $\mathsf{root}_j$
for the resulting root and
$\mathsf{Path}_{j,x}$ for the leaf's sibling hashes
and orientation bits. A \emph{public opening shell}
consists of the index $x$, the group commitment
vector $\mathsf{VC}(Q_j(x);\mathbf r_{j,x})$
and its Merkle authentication path.
It does \emph{not} reveal $Q_j(x)$ or
$\mathbf r_{j,x}$.
\end{definition}

\begin{lemma}[Perfect privacy of the authenticated shell]\label{lem:full:merkle-shell}
Fix any finite collection of field words
$w_j:\Omega_j\to\F$. Build all the Merkle
trees with independent uniform Pedersen
blindings at every $(j,x)$.
Then the joint distribution of all roots,
all commitment vectors and \emph{all} their
Merkle paths is independent of the words
$(w_j)_j$. The result holds for any
deterministic hash $H$.
\end{lemma}
\begin{proof}
By Lemma~\ref{lem:full:leaf-uniform},
the complete public array of bit-commitment
group elements is a product uniform
distribution independent of the words.
All leaf hashes, internal hashes, roots
and paths are deterministic functions of
that public array and fixed indices.
Applying the same deterministic function
to equal input distributions yields equal
output distributions.
\end{proof}

\begin{definition}[Authenticated shell opening]\label{def:full:auth-commit}
For a public root $\mathsf{root}_j$,
the authenticated shell at $(j,x)$
reveals a \emph{commitment vector}
$\mathbf C_{j,x}$ and its Merkle path,
but not its field-symbol opening. The verifier recomputes
the leaf and root and checks equality with
$\mathsf{root}_j$.
Any assertion about the hidden symbol
$a_{j,x}\in\F$ and the blinding $\mathbf r_{j,x}$
is part of a separate zero-knowledge
\emph{relation proof}, not of the
public opening shell.
\end{definition}

\begin{lemma}[Merkle opening conflict implies a hash collision]\label{lem:full:merkle-binding}
Fix one root and index, and assume canonical
injective serialization and proper domain
separation. If two valid authentication shells
for that root and index contain distinct
commitment vectors or disagree at a node
of the authenticated path, then these shells
yield two distinct inputs with the same hash
output, unless their entire authenticated
leaf-to-root sequences are identical.
\end{lemma}
\begin{proof}
Compare the two purported authenticated
chains from the root downward.
They share the same terminal root hash.
If the underlying child encodings at the
highest differing node are different,
but their parent hashes agree, the two
distinct serialized parent inputs collide
under $H$. If no internal node differs,
then at the leaf the distinct commitment
vectors give distinct injectively serialized
leaf inputs with the same leaf hash,
again a collision. Domain separation
prevents a leaf encoding from coinciding
with a node encoding.
\end{proof}

\begin{remark}[Exact authentication and privacy distinction]
Collision resistance bounds malicious
changes to authenticated commitment vectors.
Discrete-log binding bounds changes to
their hidden bit values.
Neither property makes a \emph{raw}
unmasked field-symbol opening private.
Accordingly the interface above never sends
the field symbol itself; a zero-knowledge
relation proof must enforce its opening
and all arithmetic constraints.
\end{remark}
